\chapter{Your first run}\label{ch:run}

\emph{You want an issue implemented, reviewed and merged by agents, and to stay in control while
it happens.} A pipeline run does exactly that (figure~\ref{fig:pipeline} in
chapter~\ref{ch:what}). This chapter follows one run from start to merge, and shows what to do
when a run needs you, fails or should stop.

\section{Start the run}\label{sec:start-run}

\begin{enumerate}
  \item Open the project and find the issue in the issue list.
  \item Choose \ui{Start run} on it. (In the chat, \cmd{/factory run 57} does the same, and the chat
    agent can start runs when you ask it to.) Starting needs your Claude credential; without one
    Agern answers \msg{Register a Claude credential in Settings first}.
  \item Agern opens the run page (figure~\ref{fig:run-started}). You are the run's owner: it runs
    on your credentials, and only you can approve it or talk to it.
\end{enumerate}

\screenshot{33-run-started-desktop}{A run that has just started: the spec-writer is at work.}{run-started}

\section{Read the run page}\label{sec:run-page}

The head of the page shows the issue, the run's status (for example \ui{spec}, \ui{building},
\ui{reviewing}, \ui{needs human}, \ui{merged}), the owner, the current round, the branch the
builder works on, the pull request once there is one, and how long the run has been going. The
\ui{SANDBOX} line says whether the builder uses a sandbox (chapter~\ref{ch:sandboxes}).

Below it is the pipeline diagram. Each box is a stage with its model and state:
\ui{pending}, \ui{running} with a timer, \ui{done} with its verdict, or failed. The arrows show
where the run goes next. The row of chips under the diagram lists the stages of each round.

\begin{enumerate}
  \item Watch the stages light up one after another: the spec-writer posts its specification on the
    issue, the builder implements it, commits, pushes and opens a pull request, and the four
    reviewers start together (figure~\ref{fig:run-live}).
  \item Select a box to see that stage below the diagram: its status, verdict, model, time, cost and
    turns, its summary and findings, and its full transcript, live while it runs.
  \item If any reviewer answers CHANGES, the same builder conversation gets all findings and fixes
    them in a fix round; then the four review again. The number of fix rounds is limited
    (two by default, see section~\ref{sec:models}); a selector above the stage details switches
    between rounds.
  \item When all four pass, the approver decides. If it approves, Agern squash-merges the pull
    request and closes the issue.
\end{enumerate}

\screenshot{50-run-reviewing-desktop}{A run in review: the four reviewers work in parallel.}{run-live}

On a phone the run page is one column, with the diagram on top (figure~\ref{fig:run-phone}).

\phoneshot{52-run-phone}{A run on a phone.}{run-phone}

\section{Reviewer screenshots}\label{sec:reviewer-screenshots}

The UX and UI design reviewers do not only read the change: they run the app under review in a
headless browser and attach screenshots to their stage, each with a caption
(figure~\ref{fig:reviewer-screenshots}). Select a thumbnail to see it large. If a change touches
the user interface and no reviewer attached a screenshot, the run page warns
\msg{UI change reviewed without screenshots}, as in figure~\ref{fig:run-merged}. This needs the
headless browser installed in section~\ref{sec:agent-tools}.

\screenshot{32-run-reviewer-screenshots-desktop}{A UX reviewer's verdict, finding and screenshots.}{reviewer-screenshots}

\section{Steer a running agent}\label{sec:steer}

You can talk to the spec-writer or the builder while it works, for example to point it at an
existing helper or to say what to leave alone.

\begin{enumerate}
  \item Select the running spec-writer or builder in the diagram.
  \item Below its transcript, type into \ui{Tell the spec-writer something} (or the builder) and
    press Enter (figure~\ref{fig:run-steer}).
  \item Your message is queued and picked up at the agent's next step, after its current tool call.
    It is not posted on GitHub.
\end{enumerate}

\screenshot{49-run-steer-desktop}{Sending a message to the running spec-writer.}{run-steer}

\section{When a run needs a human}\label{sec:needs-human}

When the approver escalates, or a step needs a decision, the run stops in \ui{needs human}. The
page says why, in a red box under the header (figure~\ref{fig:run-needs-human}). The projects
page lists it under \ui{needs human} as well.

\begin{enumerate}
  \item Read the reason, the reviewers' findings and the pull request on GitHub.
  \item If you have questions, choose \ui{Talk to the agent}. Ask why the run needs a human, or ask
    for a change (figure~\ref{fig:run-talk}). A change is committed, pushed and reviewed again by
    the four reviewers before the run can be approved. Talking never approves anything, and the
    conversation is not posted on GitHub.
  \item When you are satisfied, choose \ui{Approve \& merge}. If the approver advised against it,
    Agern asks once more (\ui{Merge against the approver's advice?}); \ui{Merge anyway} squash-merges
    the pull request, \ui{Don't merge} keeps waiting. \cmd{/factory approve} in the chat does the same.
\end{enumerate}

\screenshot{51-run-needs-human-desktop}{A run waiting for its owner after the approver escalated.}{run-needs-human}

\screenshot{34-run-talk-desktop}{Talking to the agent about a run that needs a human.}{run-talk}

\section{Cancel a run}\label{sec:cancel}

\begin{enumerate}
  \item Choose \ui{Cancel run} in the run's header. Only the owner and admins can.
  \item Confirm with \ui{Cancel run}; \ui{Keep running} changes your mind. The agents stop; the branch
    and any pull request stay on GitHub.
\end{enumerate}

\section{Retry a failed run}\label{sec:retry}

A run can fail, for example when a stage takes longer than its time limit, your usage limit is
reached, or the server restarts while it runs. The run page shows the error in red.

\begin{enumerate}
  \item Choose \ui{Retry\ldots}. The dialog offers the ways that fit and preselects the sensible one,
    with its reason (figure~\ref{fig:run-retry}):
    \begin{itemize}
      \item \ui{Update and merge}: only when the work is finished and approved, but the merge was
        refused because the target branch moved on meanwhile (a database migration with the same
        version, or a merge conflict). Agern merges the target into the run's branch, the builder
        resolves the conflict, and the reviewers and the approver run again before it merges; no new
        spec and no new build. A run that waits for you offers it as a button of its own,
        \ui{Update and merge}, before \ui{Approve \& merge}; in the chat it is
        \cmd{/factory approve <run> update}. Each run has two update rounds (an administrator can
        change that with \envvar{FACTORY_MAX_UPDATE_ROUNDS});
      \item \ui{Retry the failed step}: only when the builder finished and a later step (push,
        pull request, merge) failed;
      \item \ui{Continue from the commits}: a new run starts from the old branch, reuses the spec
        and finishes the work;
      \item \ui{Start over}: a new run on a new branch from scratch; the old branch is left as it is.
    \end{itemize}
  \item Confirm with the button, which repeats your choice.
\end{enumerate}

\screenshot{35-run-retry-desktop}{Retrying a failed run.}{run-retry}

\section{The merged result}\label{sec:merged}

A run that the approver (or you) approved ends as \ui{merged} (figure~\ref{fig:run-merged}). The
pull request is squash-merged into the default branch, the issue is closed on GitHub, and the run
page keeps the whole history: every stage, its transcript, findings and screenshots, and the cost.

\screenshot{31-run-merged-desktop}{A merged run, here with the warning that a UI change was reviewed without screenshots.}{run-merged}

\whatyouhave
You have taken an issue to a merged pull request, and you know how to follow a run, steer its
agents, decide when it needs you, and cancel or retry it. Next, plan several issues as one
release: chapter~\ref{ch:releases}.
